\section{Evaluation Procedure}
%szb: ugyanazok a gondolataim erről a section-ről, mint amiket az előbbiekben leírtam. Az egész section-ben csak ismételsz, nagyon rontja ez a thesis minőségét. Ha hosszt akarsz nö

The evaluation process was designed to compare the performance of the base and
fine-tuned models under identical conditions. All experiments use the same
translation direction, dataset split, prompt format, and evaluation metrics in
order to ensure fair comparison.

For each example in the evaluation dataset, the model receives a Java source
code snippet and is required to generate the corresponding C++ translation. The
generated output is then cleaned and post-processed before evaluation. This
post-processing step removes markdown formatting, code fences, or explanatory
text that may be generated by instruction-tuned models.

The evaluation itself is performed using two complementary approaches:
similarity-based evaluation and execution-based evaluation.

The first metric is BLEU, which measures token-level similarity between the
generated C++ code and the reference translation. BLEU is calculated for each
generated sample and then averaged across the entire evaluation dataset. Higher
BLEU scores indicate greater lexical similarity between the generated and
reference code, although BLEU does not guarantee that the translated code is
functionally correct.

The second evaluation method is execution-based testing using the
CodeEditorBench framework. In this step, the generated C++ code is compiled and
executed against predefined test cases. A translation is considered successful
only if it compiles correctly and produces the expected outputs for all tests.
This provides a more realistic measure of correctness than token-level metrics
alone.

The CodeEditorBench evaluation was executed inside a Docker container using the %szb: neeeeeeee megint a docker......... Ez ugye ai-generated?
official benchmark image. The use of Docker ensures that all generated programs
are tested under the same environment, including identical compiler versions,
runtime dependencies, and operating system settings. This improves
reproducibility and reduces the risk that differences in execution results are
caused by environment-related factors rather than model quality.

For each model, both BLEU score and execution-based metrics are collected and
compared. The base and fine-tuned variants of StarCoder2 and Qwen are evaluated
separately, making it possible to analyze whether fine-tuning leads to
improvements in syntactic similarity, compilation success, and functional
correctness.

In addition to aggregate scores, the generated outputs are also inspected
manually in selected cases. This allows the identification of common error
patterns, such as syntax errors, incorrect data types, missing library imports,
incorrect memory handling, or logical inconsistencies in the translated code.
Such qualitative analysis complements the numerical evaluation and helps explain
differences between the models.

The evaluation procedure therefore combines quantitative and qualitative %szb: az egész section egy ismétlés, kb. semmi újat nem mond, cserébe még az ismétlést is összefoglalod - sokkal egyenesebbre kellene straightforward-abbá kellene az egész method-ot alakítani
analysis. BLEU provides a fast similarity-based measure, while CodeEditorBench
assesses actual execution behavior. Combining these two approaches makes it
possible to analyze whether high similarity scores correspond to correct and
executable translations, which is one of the central research questions of this
thesis.